% RESEARCH_DRAFT_SHA256: 4b147b8b6aaa50c97553724549d2da220e7c2bc2d997f1479bbfab77b10c7c96
% Layout and source-reference wrappers expanded for the web renderer.
\documentclass[11pt]{article}
\usepackage[margin=25mm]{geometry}
\usepackage{fontspec,unicode-math}
\setmainfont{Latin Modern Roman}
\setmathfont{Latin Modern Math}
\usepackage{amsmath,amssymb,mathtools,xurl,hyperref}
\hypersetup{colorlinks=true,urlcolor=blue,linkcolor=blue}
\setcounter{secnumdepth}{0}
\setlength{\parindent}{0pt}
\setlength{\parskip}{5pt}
\setlength{\emergencystretch}{3em}
\newcommand{\sref}[2]{\hyperlink{source:#1:#2}{[S#2]}}
\newcommand{\cataloguescope}[1]{\paragraph{Source record.}#1}
\newcommand{\floor}[1]{\left\lfloor#1\right\rfloor}
\begin{document}
\section*{Research-attempt provenance}
This standalone research attempt retains the complete original definition body below,
including its historical source and preparation notes. Those notes describe the
earlier catalogue preparation; the new attempt and its verification are identified
separately after that body. No original ID, catalogue status, or source statement is changed.

% UnsolvedMath ID 30006644; source catalogue code OWR-14299914-005
\section{Convergence of Accelerated Sampling on Probability Spaces}\label{30006644}
{\small UnsolvedMath ID 30006644 · Computer Science}
\subsection{Definitions and mathematical statement}

\paragraph{Shared notation.}Standard number systems are $\mathbb N,\mathbb Z,\mathbb Q,\mathbb R,\mathbb C$. Specialized notation and index ranges follow the cited source and the definitions below.


A fixed-point iteration seeks a measure $\mu$ with $T(\mu)=\mu$. Anderson mixing combines a finite history of residuals to accelerate the iteration. Stein geometry uses the source\textquotesingle s kernel/operator structure to define descent for probability distributions.
\paragraph{Record scope and qualifications.}
The catalogue title is retained as metadata. Named certificates, contracts, bounds or frameworks in that title are not asserted as results of the original source extract.
The iteration map, mixing parameters, topology, regular subsets, and Stein metric are not specified in the supplied question.

Source-record limitation: The fixed-point map, accelerated update, regular measure subsets, and convergence topology are not supplied.
\paragraph{Statement recorded in the supplied source.}
Prove convergence of the proposed accelerated fixed-point and Anderson-mixing methods on suitable regular subsets of general probability measures, including the Stein-geometry sampling algorithm for smooth densities. \href{https://doi.org/10.4171/owr/2026/13}{[S2]}
\paragraph{Source status.}
The supplied catalogue labels this source record partially solved; that metadata was checked against the supplied snapshot on 2026-10-11. The label is preserved for provenance and is not an independent verification of present open status or mathematical correctness. The catalogue literature assessment is dated 2026-08-22.
\subsection{Short English statement}
Prove convergence of the cited accelerated probability-measure iterations beyond the Gaussian specialization.
\subsection{Sources}
\begin{itemize}
\item[S1] Ulam AI and the UnsolvedMath Contributors, UnsolvedMath: A Curated Collection of Open Mathematics Problems, record 30006644 (OWR-14299914-005). Catalogue curation and metadata: CC BY 4.0. The source extract is retained; definitions, English summaries and qualifications are editorial additions.
\url{https://huggingface.co/datasets/ulamai/UnsolvedMath}.
\item[S2] Flows on Measure Spaces and Applications in Machine Learning (2026), Oberwolfach Reports, DOI 10.4171/owr/2026/13, pp. 795--884
\url{https://doi.org/10.4171/owr/2026/13}.
\item[S3] Aksenov--Eigel--Oster, Anderson Mixing in Bures--Wasserstein Space
\url{https://arxiv.org/abs/2601.22038}.
\item[S4] Aksenov--Eigel, An Eulerian approach to regularized JKO
\url{https://doi.org/10.1007/s10915-025-03027-4}.
\item Source-specific Oberwolfach Reports licences, including CC BY-SA 4.0 where applicable, are retained; see THIRD\_PARTY\_NOTICES.md.
\end{itemize}
\paragraph{Preparation and review.}
The original source and UnsolvedMath contributors retain attribution for the source extract. Codex (GPT-6; exact serving revision not exposed) assisted with definitions, English summaries, scope notes and formatting. No new human mathematical authorship or independently verified solution is asserted. Permanent source ID 30006644 and code OWR-14299914-005 checked against the complete supplied registry and current repository. Distinct upstream records, including repeated questions and broad or damaged extracts, are retained for complete coverage of the source CS category. A repeated record is not a new independent conjecture. Complete mathematical-equivalence and current-literature review has not been performed. Preparation and final review follow contribution prompts C and D, with source gaps disclosed. Structural and rendering checks do not establish a new theorem or open status.
\label{end:30006644}

\section{Literature and context check}
The primary report, printed pp. 879--880, describes local Gaussian/Bures--Wasserstein results and explicitly leaves smooth-density Stein-geometry convergence open. The update and geometry needed for a general theorem are missing from the abbreviated record. The following exact two-point probability example tests unrestricted residual mixing, then proves a conditional safeguard lemma; neither is asserted to be the source's full proposed method.
\section{Research attempt: Anderson extrapolation can leave probability space}

Identify Bernoulli distributions with $p\in[0,1]$ and use absolute distance (equivalently total variation). Define the contraction
\[
T(p)=\frac{31}{50}+\frac9{10}\max\{p-\tfrac45,0\}.
\]
It maps $[0,1]$ into $[31/50,4/5]$, has Lipschitz constant $9/10$, and has fixed point $31/50$. Starting at $p_0=1$, ordinary iteration gives $p_1=T(p_0)=4/5$, with residuals
\[
r_0=T(p_0)-p_0=-1/5,\qquad
r_1=T(p_1)-p_1=-9/50.
\]
The unconstrained two-history Anderson step chooses coefficients summing to one to minimize the absolute mixed residual. Here its exact zero-residual choice is $\alpha_0=-9$, $\alpha_1=10$. The resulting proposed iterate is
\[
\alpha_0T(p_0)+\alpha_1T(p_1)
=-9\cdot\frac45+10\cdot\frac{31}{50}=-1.
\]
It is not a probability. Thus even genuine consecutive iterates of a contraction do not justify unconstrained affine mixing in probability coordinates. Smoothness alone does not repair the example: replace the maximum by $\epsilon\log(1+e^{(p-4/5)/\epsilon})$. For sufficiently small positive $\epsilon$ this smooth map still maps the interval into itself, has derivative between zero and $9/10$, and its two-step Anderson output remains negative by continuity of the residual formula at the displayed nonzero residual difference.

\paragraph{A conditional valid-space safeguard.}
Let $(X,d)$ be complete and $T:X\to X$ be a contraction of factor $0\leq q<1$. Its unique fixed point $x_*$ exists: the ordinary iterates are Cauchy because consecutive distances are bounded by a geometric sequence, their limit is fixed by continuity, and two fixed points would satisfy $d\leq qd$. Define residual $R(x)=d(x,Tx)$. The triangle inequality gives
\[
d(x,x_*)\leq R(x)+q d(x,x_*),\qquad
R(Tx)\leq q R(x).
\]
Fix $0\leq\theta<1$. At each step accept any accelerated candidate $y\in X$ only if $R(y)\leq\theta R(x)$; otherwise use $Tx$. Then with $\rho=\max(q,\theta)<1$,
\[
R(x_k)\leq\rho^kR(x_0),\qquad
d(x_k,x_*)\leq\frac{\rho^kR(x_0)}{1-q}.
\]
This follows by induction and the first residual bound. Membership in $X$ must be checked before acceptance; the negative Bernoulli proposal fails it. This is a complete convergence proof under a specified contraction and computable exact residual, but it neither promises acceleration nor proves that a Stein sampling map satisfies those hypotheses.
\section{Verification and adversarial audit}
The residual-minimizing coefficients and negative output are exact rational numbers and arise from consecutive unaccelerated iterates. Smooth approximation preserves a nonzero residual denominator and hence the sign for small enough smoothing. The safeguard requires completeness, a self-map, a strict contraction, and actual residual tests; dropping any of these is not justified by the lemma.
\paragraph{Audit corrections.} Used an actual consecutive-history counterexample, rather than arbitrary unrelated iterates, and distinguished residual decrease from preservation of probability measures.
\paragraph{External computation.} None; symbolic arguments and boundary checks only.
\section{Scope and remaining gaps}
The source's regular density subsets, Stein metric, retractions, transport, sampling errors, and contraction/alternative convergence conditions are not established. A conditional safeguarded contraction result does not resolve the general accelerated sampling agenda.
\paragraph{Integration review.} Codex (GPT-6; exact serving revision not exposed) reviewed the written arguments and expanded layout and source-reference wrappers. Repository metadata and the production web renderer were checked using Python 3.13 and JavaScript with MathJax 3.2.2. This model-assisted review is not independent human proof verification; no PDF compilation or formal proof checking is asserted.

\section{Completion Estimate}
\noindent\textbf{COMPLETION: 0\%}
The calculations below the preserved source statement investigate a concrete route,
but establish no substantiated new progress toward the entire original target.
Known special cases and failed sufficient conditions are not credited as a solution.
\section{References}
The original source references and catalogue attribution are preserved in the
Sources section above. The following URLs were actually consulted for this attempt;
retention of an additional original reference does not claim fresh consultation of it.
\begin{itemize}
\item \url{https://ems.press/content/serial-article-files/53606} (consulted 2026-10-11). 
\end{itemize}
\section{Model and contributor disclosure}
Codex GPT-6 completed written-proof review of this attempt at its declared scope; this is model-assisted review, not independent human proof checking.
OpenAI gpt-6-astra, requested reasoning configuration ultra, generated this research
attempt on 2026-10-11. The exact serving revision is not exposed. Codex GPT-6
performed orchestration and repository-format review; its exact serving revision
is also not exposed. The mathematical route, boundary checks and final adversarial
self-audit followed contribution prompts A then B in
\texttt{docs/contribution\_prompts.md}, after reading \texttt{CONTRIBUTING.md}.
Python 3.14.2 assembled the standalone documents; any mathematical computation is
specified in Verification. No theorem prover was used. No human mathematical
authorship, substantial human proof correction, or independent human proof checking
is asserted. Model-assisted review, source retrieval, and structural checks do not
establish correctness or novelty. Full private reasoning transcripts are unavailable
as reproducibility artifacts; the complete arguments, inputs and outcomes actually
used are stated here or in the companion batch audit.
\section{Final status}
Unresolved as a new research contribution. This is a concrete route investigation,
not a claim of new complete resolution or a current-status certification. A known
answer reproduced for context does not become a new contribution. The catalogue's
source-status label is preserved; this attempt status does not assert that a
source-labelled solved question is open.
\end{document}
